% \section{Conclusion}
% This paper provides an analysis of the determinants of aggregate market elasticity in a general equilibrium framework with rich investor heterogeneity, passive investment, and financial constraints. Our analysis yields several important insights about market inelasticity and its implications for asset prices and market volatility.

% A central finding of our work is the crucial role of cross-price elasticity--the sensitivity of demand for risky assets to changes in interest rates. When cross-elasticity is zero, aggregate market elasticity is simply an average of individual investors' price elasticities, consistent with partial equilibrium intuition. However, with non-zero cross-elasticity, the market can be infinitely elastic even when individual investors are highly inelastic, as changes in interest rates offset movements in risk premia. This highlights how general equilibrium effects fundamentally alter the relationship between individual and aggregate elasticities.

% We show that beyond cross-elasticity effects, the key determinant of market elasticity is how risk is allocated in the economy. When passive investors hold an efficient share of risky assets, the market remains infinitely elastic regardless of individual investor preferences. However, when risk is misallocated, portfolio flows have meaningful price impacts.

% Our model demonstrates that market elasticity is both state-dependent and time-varying, influenced by the distribution of wealth between active and passive investors as well as the allocation of wealth among active investors themselves. While passive investment and leverage constraints amplify the price impact of flows, preference heterogeneity can make markets more elastic by improving risk allocation. These findings highlight that market inelasticity serves as an indicator of underlying inefficiencies in risk allocation rather than a fundamental feature of markets.

% The framework developed here helps explain several empirical patterns, including the countercyclical nature of the volatility multiplier and the relationship between passive investment growth and market dynamics. It also provides a foundation for analyzing how different market frictions interact to influence aggregate elasticity and market stability. Our results suggest that policies aimed at improving risk allocation may be more effective at reducing excess volatility than those focused on individual investor behavior.

% Future research could extend this framework to study the implications of market inelasticity for asset pricing anomalies, monetary policy transmission, and financial stability. The role of market structure and trading mechanisms in determining aggregate elasticity also remains an important area for investigation.

\section{Conclusion}

This paper develops a general equilibrium model with heterogeneous investors, passive demand, and financial constraints to study the determinants of aggregate market elasticity. We show how general equilibrium adjustments, particularly the joint movement of the risk-free rate and the risk premium, fundamentally shape the price impact of portfolio flows.  

A central insight is the role of cross-price elasticity. When demand for risky assets responds to interest rates, shifts in the risk premium are offset by changes in the risk-free rate. In this case, the market can be infinitely elastic even if individual investors are highly inelastic. Aggregate elasticity becomes finite only when risk is initially misallocated: portfolio flows then alter aggregate savings behavior, preventing interest rates from fully offsetting risk-premium movements.  

Macro elasticity is both state-dependent and time-varying. Passive investors amplify price impacts, leverage constraints further tighten risk-bearing capacity, while preference heterogeneity increases elasticity by improving risk allocation. Inelasticity alone does not generate excess volatility, nor do flows in infinitely elastic markets; both elements are needed to explain observed fluctuations and the countercyclical volatility multiplier.  

Our analytical results rely on a state-global perturbation method that provides closed-form characterizations of elasticity, while our quantitative analysis solves the full dynamic model numerically. The calibrated model highlights how the wealth distribution across households, intermediaries, and long-term investors shapes aggregate elasticity and volatility.  

Overall, our findings suggest that market inelasticity is best understood as a symptom of inefficient risk allocation, rather than a structural feature of financial markets. Improving the allocation of risk across sectors may therefore be more effective for reducing excess volatility than targeting individual investor behavior. Future work could extend this framework to study asset pricing anomalies, monetary policy transmission, and the role of market structure in shaping aggregate elasticity.
